% !TeX root = ../main.tex
% !TeX encoding = utf8
	\section{GIỚI THIỆU TỔNG QUAN VÀ BỐI CẢNH NGHIÊN CỨU}
	
	\subsection{Bối cảnh và động lực nghiên cứu}
	Trong lý thuyết suy diễn thống kê truyền thống, việc xây dựng khoảng tin cậy cho một tham số tổng thể $\theta$ (như trung bình $\mu$, phương sai $\sigma^2$, hay tỷ số Odds Ratio) thường dựa vào các giả định phân phối tham số chặt chẽ (chẳng hạn như giả định mẫu rút ra từ phân phối Chuẩn $N(\mu, \sigma^2)$) hoặc dựa trên các kết quả tiệm cận khi kích thước mẫu $n \to \infty$ thông qua Định lý Giới hạn Trung tâm (CLT) \cite{buithuytrang2026}. 
	
	Trong một số ứng dụng, dữ liệu có thể không đáp ứng tốt các giả định này; một số tình huống thường gặp gồm \cite{efron1993}:
	\begin{itemize}
		\item \textbf{Cỡ mẫu nhỏ:} Chi phí, độ hiếm của đối tượng hoặc thời gian thu thập có thể giới hạn cỡ mẫu. Khi đó, xấp xỉ tiệm cận có thể chưa chính xác; mức độ ảnh hưởng phụ thuộc vào bài toán cụ thể \cite{efron1993}.
		\item \textbf{Phân phối lệch hoặc có ngoại lai:} Thu nhập, chi phí bồi thường bảo hiểm, thời gian sống sót và nồng độ độc tố có thể có phân phối lệch hoặc đuôi dày \cite{davison1997}.
		\item \textbf{Thống kê ước lượng phức tạp:} Trong nhiều bài toán thực tế, nhà nghiên cứu quan tâm đến các tham số không có công thức phương sai giải tích đơn giản, chẳng hạn như Trung vị (Median), Hệ số Gini, Tỷ số rủi ro (Hazard Ratio), hay các trọng số tối ưu trong danh mục đầu tư \cite{efron1993}.
	\end{itemize}
	
	Khi các giả định của phương pháp cổ điển không phù hợp, khoảng thu được có thể phản ánh chưa đúng độ bất định và coverage thực tế có thể khác mức danh nghĩa $1-\alpha$ \cite{davison1997}. Bootstrap là một phương pháp khác để ước lượng độ bất định, nhưng hiệu năng của khoảng vẫn phụ thuộc vào dữ liệu và cách lấy mẫu lại \cite{efron1979}.

	\subsection{Mục tiêu và phạm vi nghiên cứu}
	Báo cáo tập trung vào hai phương pháp khoảng tin cậy Bootstrap: \textbf{Normal CI} và \textbf{Basic CI} \cite{buithuytrang2026}. Percentile CI được đưa vào phần mở rộng.
	
	Báo cáo hướng đến các mục tiêu sau:
	\begin{enumerate}
		\item \textbf{Trình bày cơ sở toán học:} Giới thiệu ECDF, một số tính chất cơ bản, nguyên lý Plug-in và mô hình xác suất Bootstrap \cite{efron1993}.
		\item \textbf{So sánh phương pháp:} Trình bày công thức Normal CI và Basic CI, đồng thời giải thích phép biến đổi phân vị trong Basic CI \cite{davison1997}.
		\item \textbf{Minh họa và mô phỏng:} Phân tích hai ví dụ bằng cách liệt kê chính xác trong R; mô phỏng Monte Carlo với $K=10{,}000$ lần lặp để ước lượng coverage và độ dài trung bình trong bốn kịch bản.
		\item \textbf{Ứng dụng và mở rộng:} Mô tả một ứng dụng BCa trong bài báo khoa học, so sánh Percentile CI và khảo sát tính bất biến dưới phép biến đổi $g(\theta)$; khuyến nghị giới hạn theo kết quả đã trình bày \cite{deng2013, efron1993}.
	\end{enumerate}
	
	Phạm vi lý thuyết tập trung vào dữ liệu i.i.d. và ba phương pháp Normal, Basic, Percentile CI. Phần mô phỏng chỉ xét hai phân phối, hai cỡ mẫu và các thiết lập nêu tại Phần 6; kết quả không đại diện cho mọi dạng dữ liệu \cite{davison1997}.
	
	\subsection{Bối cảnh ra đời của thống kê phi tham số và resampling}
	Sự phát triển của một số phương pháp lấy mẫu lại gắn với tiến bộ về tính toán. Các mốc tiêu biểu gồm:
	\begin{itemize}
		\item \textbf{Kiểm định hoán vị (thập niên 1930):} Ronald Fisher (1935) và E. J. G. Pitman (1937) phát triển các phép kiểm định dựa trên việc xét các cách gán nhãn lại dữ liệu \cite{efron1993}. Khi số cấu hình tăng theo giai thừa của cỡ mẫu, việc liệt kê đầy đủ bị giới hạn bởi chi phí tính toán \cite{efron1979}.
		\item \textbf{Jackknife (1949, 1958):} Maurice Quenouille (1949) đề xuất loại lần lượt từng quan sát để giảm độ chệch của ước lượng. John Tukey (1958) mở rộng phương pháp này để ước lượng phương sai và đặt tên là Jackknife \cite{efron1979}. Phương pháp có hạn chế với một số thống kê không trơn, chẳng hạn trung vị \cite{efron1993}.
		\item \textbf{Bootstrap (1979):} Efron giới thiệu phương pháp trong bài báo \textit{"Bootstrap Methods: Another Look at the Jackknife"} \cite{efron1979}. Bootstrap thay phân phối tổng thể bằng phân phối thực nghiệm $F_n^*$ rồi lấy mẫu lại có hoàn lại. Khi không liệt kê toàn bộ các mẫu, mô phỏng Monte Carlo được dùng để xấp xỉ phân phối Bootstrap \cite{efron1993}.
	\end{itemize}
	
	
	% ==================== PHẦN 2 ====================
